% !TEX root = /home/eb/.vscode/extensions/qftrules.programme-colle/tmp/Exercice.tex
%%%%%%%%%%%%%%%%%%%%%%%%%%%%%%%%%%%
\begin{exo}[PC][1][colle]{Autour du laser Hélium-Néon}
  \begin{quest}
    \item Citer les différents constituants d'un laser à gaz.
    \item Qu'appelle-t-on émission stimulée? Quelle est sa caractéristique principale? Citer les trois types de transitions possibles entre deux niveaux d'énergie $\varepsilon_1$ et $\varepsilon_2\left(\varepsilon_2>\varepsilon_1\right)$.
    \item Un gaz d'atomes en équilibre avec un thermostat de température $T$, obéit à la loi de Boltzmann : la population $N$ d'un niveau d'énergie atomique $\varepsilon$ s'écrit
    $$
N=\mathrm{C}^{\text {te }} \mathrm{e}^{-\varepsilon / k_{\mathrm{B}} T} \quad \text { avec } \quad k_{\mathrm{B}}=1,38 \times 10^{-23} \mathrm{JK}^{-1}
$$

Calculer le rapport $N_2 / N_1$ des populations des niveaux $\varepsilon_2$ et $\varepsilon_1$ correspondant aux niveaux $5 \mathrm{~s}$ et $3 \mathrm{p}$ du néon. On donne $\varepsilon_2-\varepsilon_1=1,96 \mathrm{eV}$ et $T=300 \mathrm{~K}$.
\item En général, une onde se propageant suivant une direction $\mathrm{O} z$ dans un milieu à deux niveaux d'énergie s'amortit exponentiellement. Son intensité suit la loi
$$
I=I_0 \mathrm{e}^{-\alpha \mathrm{z}} \quad \text { avec } \quad \alpha \text { proportionnel à }\left(N_1-N_2\right)
$$

Expliquer comment, dans un laser à gaz, le milieu gazeux peut produire une amplification de l'intensité de l'onde.
\item Dans un laser Hélium-Néon, l'émission correspond à la transition entre les niveaux $5 \mathrm{~s}$ et $3 \mathrm{p}$ des atomes de néon. Calculer la fréquence laser $v_0$ et la longueur d'onde correspondante $\lambda_0$. Dans quel domaine spectral ce laser émet-il?

On considère un laser Hélium-Néon dans lequel la cavité est constituée de deux miroirs plans $\mathbf{M}$ et $\mathbf{M}^{\prime}$ séparés de la distance $L=20,0 \mathrm{~cm}$. Le miroir $\mathrm{M}$ est parfaitement réfléchissant alors que $\mathrm{M}^{\prime}$ présente un pouvoir de réflexion $R=99,6 \%$. Lindice de réfraction du milieu amplificateur est égal à 1 .
\item Énoncer la condition de résonance de la cavité. En déduire les fréquences des modes propres de la cavité ainsi que l'intervalle spectral libre $\Delta v$ correspondant à l'écart fréquentiel entre deux modes propres successifs.
\item En notant $\tau$ le temps mis par la lumière pour faire un aller-retour dans la cavité, on peut montrer que la largeur spectral d'un mode laser vaut
$$
\delta v_{1 / 2} \simeq \frac{|\ln R|}{2 \pi \tau}
$$

Calculer $\delta v_{1 / 2}$ et l'écart $\delta \lambda_{1 / 2}$ correspondant.
\item Un oscillateur excité sur l'une de ses fréquences propres présente une réponse résonante caractérisée par un nombre sans dimension $Q$, appelé facteur de qualité. Il vaut
$$
Q=\frac{v_0}{\delta v_{1 / 2}}
$$

Calculer le facteur de qualité de ce laser Hélium-Néon. Quelle propriété du laser est mise en évidence ici?
\end{quest}
\end{exo}
%%%%%%%%%%%%%%%%%%%%%%%%%%%%%%%%%%%

%%%%%%%%%%%%%%%%%%%%%%%%%%%%%%%%%%%
\begin{exo}[PC][1][oral]{Propriétés d'une diode laser}
  Déterminer les amplitudes des champs électrique et magnétique créés par la diode laser suivante, ainsi que les caractéristiques du faisceau gaussien produit.\\[0.3cm]

Diode laser rouge Classe II @635 nm version sur tige bas profil Ref : 203324.
Caractéristiques techniques :\\[0.3cm]
\begin{minipage}{0.49\textwidth}
  \begin{liste}
  \item Puissance : 1 mW (classe 2)
  \item Extrémité : filetée M20
  % \item Maintien : tige D 10mm - L $130 \mathrm{~mm}$
  % \item Dimensions tube : diamètre $25 \mathrm{~mm}$ / longueur $90 \mathrm{~mm}$
  \item Longueur tube : $\SI{90}{mm}$
  \item Longueur d'onde : $\SI{635}{nm}$
  \item Type : diode laser
  \item Diamètre faisceau à $\SI{5}{m}$ : $\SI{9}{mm}$
  \end{liste}
\end{minipage}\hfill
\begin{minipage}{0.49\textwidth}
  \begin{liste}
  \item Divergence : $\SI{0.9}{mrad}$
  \item Diamètre faisceau en sortie : $\SI{1}{mm}$
  \item Température d'utilisation : 10 à $40^{\circ} \mathrm{C}$
  \item Polarisation : linéaire
  \item Transformateur (fourni) : $6$-$9$V DC
  \end{liste}
\end{minipage}
\sol{
\begin{liste}
  \item Puissance moyenne surfacique pour une onde sinusoïdale donnée par la moyenne du vecteur de Poynting $\moy{\Pi} = \frac{1}{2}\epz E_0^2$, d'où la puissance $P = \frac{1}{2}\epz E_0^2 \frac{\pi d^2}{4} = \SI{1}{mW}$. L'amplitude du champ électrique
\eqbox{E_0 = \sqrt{\frac{8P}{\pi\epz d^2}} = \SI{1.7e7}{V.m^{-1}}.}

  \item
  Angle de divergence $\theta = \frac{\lambda}{\pi w_0} = \SI{0.9}{mrad}$, donc col $w_0 = \frac{\lambda}{\pi \theta} = \SI{0.2}{mm}$. Longueur de Rayleigh $z_R = \frac{\pi w_0^2}{\lambda} = \SI{25}{cm}$. Le diamètre du faisceau à $D=\SI{5}{m}$ est $d \simeq 2D\theta = \SI{9}{mm}$
\end{liste}
}
\end{exo}
%%%%%%%%%%%%%%%%%%%%%%%%%%%%%%%%%%%

%%%%%%%%%%%%%%%%%%%%%%%%%%%%%%%%%%%%%%%%%%%%%%%%%%%%%%%%%%%%%%%%%%%
\begin{exo}[PC][2][TD]{Laser à rubis}
  \Notions{\item Coefficients d'Einstein \task Interaction lumière-matière \task Trois niveaux}
  Le laser à rubis fut le premier laser mis au point, conçu par T. \textsc{Maiman} en 1960. Il s’agit d’un laser à impulsion dont le milieu amplificateur est constitué de rubis, qui est un cristal d’alumine Al$_2$O$_3$ dopé aux ions chrome Cr$^{3+}$ ; ces ions sont responsables de la couleur rouge caractéristique du rubis. 
  
  L'émission laser fait intervenir $3$ états des ions Cr$^{3+}$, d'énergies $E_1<E_2<E_3$, représentés sur la \Fig{laser_rubis}. Le niveau fondamental $1$ et le niveau métastable $2$ sont responsables de l’émission laser, et associés aux coefficients d’Einstein $A = \SI{1.0e3}{s^{-1}}$ pour l'émission spontanée et $B$ pour les processus stimulés. Le niveau excité $3$ est peuplé par pompage optique au moyen de flashs lumineux, de taux $\Gamma$ (en $s^{-1})$, et se désexcite spontanément de façon non radiative, avec une durée de vie moyenne $\tau = \SI{50}{ns}$. On note $u(\omega)$ la densité d'énergie volumique du rayonnement à la pulsation laser $\omega$, de longueur d’onde $\lambda = \SI{694}{nm}$.

  \begin{fig}[laser_rubis]
    \centering
    \includegraphics[scale=1.2]{laser_rubis.pdf}
    % \capt{Diagramme énergétique des ions Cr$^{3+}$.}
  \end{fig}

% Les populations des atomes sont notées respectivement $n_1, n_2$ et $n_3$. On réalise un pompage optique entre les niveaux $E_1$ et $E_3$ à l’aide de flashs électriques. Le nombre d'atomes excités par unité de temps par le pompage optique est noté $\Gamma n_1$. On suppose que la désexcitation du niveau 3 vers le niveau 2 est non radiative (vibration du cristal) et que le temps de vie moyen dans le niveau 3 vaut $\tau \simeq \SI{e-8}{s}$. L'émission laser correspond à la transition du niveau 2 vers le niveau 1.
% \begin{fig}[]
%   \centering
%   \includegraphics[width=0.5\textwidth]{laser3.png}
%   % \capt{}
% \end{fig}
%  On utilisera les coefficients d'Einstein $A$ et $B$ pour rendre compte des phénomènes d'émission spontanée d’une part, et d’émission stimulée et d'absorption d’autre part. 

 % donnée par la loi de Planck.

\begin{quest}
\item Établir l'équation d'évolution des populations atomiques $n_1(t)$, $n_2(t)$ et $n_3(t)$, i.e. du nombre d'atomes de chrome par unité de volume dans les états $1$, $2$ et $3$ respectivement.
\sol{
$$
\dd{n_1}{t} = -\Gamma n_1 + B u(\omega) (n_2-n_1) + A n_2,
$$
$$
\dd{n_2}{t} = \frac{n_3}{\tau} - B u(\omega) (n_2-n_1) - A n_2,
$$
$$
\dd{n_3}{t} = \Gamma n_1 - \frac{n_3}{\tau}
$$
}
\item Montrer que la quantité $n_{\rm tot} = n_1+n_2+n_3$ est une constante. Estimer $n_{\rm tot}$ sachant que le rubis est constitué à $\num{0.05}\%$ en masse d’ions Cr$^{3+}$, et que la masse volumique du rubis est $\rho = \SI{4.0}{g.cm^{-3}}$.
\sol{
En sommant les trois équations précédentes, on trouve que
\eqbox{\dd{n_1}{t}  + \dd{n_2}{t}  + \dd{n_3}{t}  = \dd{n_{\rm tot}}{t} = 0.}

% Connaissant les masses molaires des constituants du rubis, on trouve que la masse molaire du rubis vaut $M \simeq M({\rm Al}_2{\rm O}_3) \simeq \SI{102}{g.mol^{-1}}$. Le nombre d'atomes de rubis par unité de volume vaut donc 
% $$n_{\rm Al_2O_3} = \frac{\rho}{M} N_A \simeq \SI{2.4e22}{cm^{-3}}.$$ Le nombre d'ions Cr$^{3+}$ par unité de volume est donc 
% \eqbox{
  % n_{\rm tot} \simeq \num{5e-5} n_{\rm Al_2O_3} \simeq \SI{1.2e19}{cm^{-3}}.
  % }
  La masse volumique d’ions Cr$^{3+}$ dans le rubis vaut $\rho_{\rm Cr} = \num{0.0005} \rho = \SI{2.0e-3}{g.cm^{-3}}$. La masse molaire de Cr$^{3+}$ est $M_{\rm Cr} = \SI{52}{g.mol^{-1}}$. Le nombre d'ions Cr$^{3+}$ par unité de volume vaut donc
  \eqbox{
    n_{\rm tot} = \frac{\rho_{\rm Cr}}{M_{\rm Cr}} N_A \simeq \SI{2.3e19}{cm^{-3}}.
  }


}
\item On pose $\Delta n=n_2-n_1$ l’écart de population. En régime stationnaire, exprimer $\Delta n$ en fonction de $n_3$. En déduire une condition nécessaire pour réaliser l'inversion de population.
\sol{
En régime stationnaire, on peut écrire un système d'équations linéaires en $n_1$, $n_2$ et $n_3$ :
$$
\begin{cases}
  0 = -\Gamma n_1 + B u(\omega) (n_2-n_1) + A n_2 , \\
  0 = \frac{n_3}{\tau} - B u(\omega) (n_2-n_1) - A n_2, \\
  0  = \Gamma n_1 - \frac{n_3}{\tau}.
\end{cases}
$$
On en déduit que $n_1 = \frac{n_3}{\Gamma \tau}$ puis la 1ere équation donne 
$$
0 = -\Gamma n_1 + B u(\omega) (n_2-n_1) + A n_2 = 0 = (A-\Gamma) n_1 + (B u(\omega) + A) \Delta n.
$$
 Ainsi,
 \eqbox{\Delta n = \frac{\Gamma - A}{A+Bu(\omega) }n_1 = \frac{\Gamma - A}{A+Bu(\omega) } \frac{n_3}{\Gamma \tau}}
% \eqbox{\Delta n = n_2-n_1 =\frac{n_3}{\tau}\left( \frac{1}{B u(\omega) + A} - \frac{1}{\Gamma}  \right).}
% \sol{
% }
% \item 
On veut $\Delta n >0$, soit $\Gamma > A$. Or $\Gamma$ est le taux de pompage et $A$ le taux d'émission spontanée. Il faut donc pomper plus vite que les atomes n'émettent aléatoirement pour inverser la population. On voit que aussi que plus les transitions non radiatives sont rapides ($\tau$ petit), meilleure est l'inversion de population.
}
\item Estimer la puissance moyenne nécessaire $P_{\rm pomp}$ fourni par pompage optique pour atteindre l'inversion de population, sachant que le volume du rubis est $V = \SI{1}{cm^3}$ et que $E_3 - E_1 = \SI{2.3}{eV}$.
\sol{
  À la limite de l'inversion de population, on a $\Gamma = A = \SI{1.0e3}{s^{-1}}$. Le nombre d'atomes excités par unité de temps est donc $A n_1 V$. On voit que dans ce cas, $\Gamma \tau \ll 1$, et donc $n_3 \ll n_1$. Comme $n_2=n_1$ à la limite de l’inversion de population, on trouve $n_1 \simeq \frac{n_{\rm tot}}{2}$. La puissance moyenne nécessaire pour atteindre l'inversion de population vaut donc
  \eqbox{
    P_{\rm pomp} = \frac{A n_{\rm tot} V}{2} (E_3 - E_1) \simeq \SI{2.2}{kW}.
  }
}
\end{quest}

\textsc{Données} : Masses molaires : $M(Al) = \SI{27}{g.mol^{-1}}$, $M(O) = \SI{16}{g.mol^{-1}}$, $M(Cr) = \SI{52}{g.mol^{-1}}$. Constante d'Avogadro : $N_A = \SI{6.02e23}{mol^{-1}}$.
\end{exo}
%%%%%%%%%%%%%%%%%%%%%%%%%%%%%%%%%%%%%%%%%%%%%%%%%%%%%%%%%%%%%%%%%%%


%%%%%%%%%%%%%%%%%%%%%%%%%%%%%%%%%%%%%%%%%%%%%%%%%%%%%%%%%%%%%%%%%%%
\begin{exo}[PC][2][TD]{Laser à trois niveaux}
  \Notions{\item Coefficients d'Einstein \task \task Pompage optique \task Trois niveaux}
  Le laser hélium-néon émet dans le rouge à la longueur d’onde $\lambda = \SI{632.8}{nm}$. 
    Les atomes de néon possèdent trois niveaux d'énergies : un niveau fondamental d’énergie $E_0$, et deux niveaux excités d’énergies $E_1$ et $E_2$, avec $E_0<E_1<E_2$. Les populations atomiques sont notées respectivement $n_0, n_1$ et $n_2$.  L'émission laser provient de la transition entre les niveaux 1 et $2$.
    
    On réalise un pompage optique entre les niveaux $0$ et $2$ en excitant les atomes d’hélium par des décharges électriques. Les atomes d’hélium transfèrent ensuite leur énergie aux atomes de néon par collision. On note $\Gamma$ le taux de pompage, et on assimile les populations de He et Ne aux niveaux $0$ et $2$.
    
    On suppose que la désexcitation du niveau 2 vers le niveau 1 donne l’émission laser, tandis que la transition du niveau 1 au niveau 0 est non radiative (collision). Le temps de vie moyen dans le niveau 1 est noté $\tau$. L'émission laser correspond à la transition du niveau 1 vers le niveau 0.
    
    On utilisera les coefficients d'Einstein $A$ et $B$ pour rendre compte des phénomènes d'émission spontanée d’une part, et d’émission stimulée et d'absorption d’autre part, entre les niveaux $1$ et $2$. On note $u(\omega)$ la densité volumique d'énergie  du rayonnement monochromatique à la pulsation laser $\omega$.
    % donnée par la loi de Planck.
    \begin{fig}[label]
      \centering
      \includegraphics[scale=1]{laser_HeNe.pdf}
      %\capt{}
    \end{fig}

\begin{quest}
\item Établir l'équation d'évolution de chaque niveau.
\sol{
  $$
  \left\{  
  \begin{aligned}
    \dd{n_0}{t} & = -\Gamma n_0 + \frac{n_1}{\tau}, \\
    \dd{n_1}{t} & = -\frac{n_1}{\tau} + B u(\omega) (n_2-n_1) + A n_2, \\
    \dd{n_2}{t} & = \Gamma n_0 -B u(\omega) (n_2-n_1) - A n_2.
  \end{aligned}
  \right.
  $$

}
\item Montrer que $n_1+n_2+n_3$ est une constante, notée $n$.
\sol{
En sommant les trois équations précédentes, on trouve que
\eqbox{\dd{n_0}{t}  + \dd{n_1}{t}  + \dd{n_2}{t}  = \dd{n}{t} = 0.}
}
\item  On pose $\Delta n=n_2-n_1$. En régime stationnaire, exprimer $\Delta n$ en fonction de $n_0$.
\sol{
En régime stationnaire, on peut écrire un système d'équations linéaires en $n_1$, $n_2$ et $n_3$ :
$$
\begin{cases}
  0 = -\Gamma n_0 + \frac{n_1}{\tau}, \\
  0 = -\frac{n_1}{\tau} + B u(\omega) (n_2-n_1) + A n_2, \\
  0  = \Gamma n_0 -B u(\omega) (n_2-n_1) - A n_2.
\end{cases}
$$
On en déduit que $n_1 = n_0 \Gamma \tau$ puis la 3e équation donne 
$$
0 = \Gamma n_0 - B u(\omega) (n_2-n_1) - A n_2 = 0 = \Gamma n_0 - A n_1 - [A+B u(\omega) ] \Delta n.
$$
 Ainsi,
 \eqbox{\Delta n = \frac{\Gamma(1 - A\tau)}{A+Bu(\omega)}n_0.}
% \eqbox{\Delta n = n_2-n_1 =\frac{n_3}{\tau}\left( \frac{1}{B u(\omega) + A} - \frac{1}{\Gamma}  \right).}
}
\item En déduire une condition nécessaire pour réaliser l'inversion de population.
\sol{
  La grandeur $A+B u(\omega)$ représente le taux d'émission total depuis le niveau $2$. On veut $\Delta n >0$, soit $$A \tau < 1.$$ La désexcitation mon radiative vers le niveau fondamental doit donc être plus rapide que l'émission spontanée vers le niveau $1$.
  Ici il n'y a pas d'effet de seuil sur le taux de pompage optique $\Gamma$, contrairement au laser au rubi. L'inversion de population est proportionnelle au pompage :  plus $\Gamma$ est grand, plus l'inversion de population est importante.

}
\end{quest}
\end{exo}
%%%%%%%%%%%%%%%%%%%%%%%%%%%%%%%%%%%%%%%%%%%%%%%%%%%%%%%%%%%%%%%%%%%

%%%%%%%%%%%%%%%%%%%%%%%%%%%%%%%%%%%%%%%%%%%%%%%%%%%%%%%%%%%%%%%%%%%
\begin{exo}[PC][2][TD]{LIDAR}
  \Notions{\item Faisceau gaussien}
Le LIDAR (\textit{Light Detection and Ranging}) est un instrument de mesure de distance par l'analyse des propriétés d'un rayonnement laser. Cette technique permet de sonder l'atmosphère sur quelques kilomètres d'altitude.
\begin{quest}
\item On envisage d'utiliser un laser YAG de longueur d'onde $\lambda_0=\SI{1,06}{\micro\meter}$ et de waist $w_0=\SI{0,40}{mm}$. Situer ce rayonnement dans le spectre électromagnétique.
\item Déterminer la divergence angulaire et la longueur de Rayleigh de ce faisceau. Déterminer la largeur du faisceau à \SI{1500}{m} d'altitude en négligeant tout phénomène d'absorption et de diffusion atmosphérique. Expliquer pourquoi il faut collimater le faisceau pour l'application recherchée.
\item Ce faisceau arrive sur un système afocal constitué d'une lentille divergente de focale $-\SI{20}{mm}$ placée à \SI{150}{mm} du waist du faisceau incident, et d'une lentille convergente de focale \SI{200}{mm}.
\begin{enumerate}
\item Faire un schéma du dispositif optique en faisant apparaître les foyers objet et image de chaque lentille. Quelle distance sépare les deux lentilles ? Pourquoi la première lentille est-elle choisie divergente ?
\item Caractériser le faisceau gaussien émergent (waist, longueur de Rayleigh, divergence angulaire).
\item Calculer la largeur du faisceau à \SI{1500}{m} d'altitude et commenter.
\end{enumerate}
\end{quest}
\end{exo}
%%%%%%%%%%%%%%%%%%%%%%%%%%%%%%%%%%%%%%%%%%%%%%%%%%%%%%%%%%%%%%%%%%

%%%%%%%%%%%%%%%%%%%%%%%%%%%%%%%%%%%%%%%%%%%%%%%%%%%%%%%%%%%%%%%%%%%%%%%%%%%%%
\begin{exo}[PC][1][TD]{Phénoménologie des pertes dans un laser}
  \Notions{\item Coefficients d'Einstein \task Pompage optique \task Deux niveaux \task Pertes}
Un laser est modélisé par un système à deux niveaux atomiques d'énergies $E_1$ et $E_2>E_1$. Soit $n_1$ et $n_2$ les densités volumiques d'atomes dans chaque niveau et $\vphi$ le flux surfacique lumineux, i.e. la densité de courant en photons. Les probabilités d'absorption et d'émission stimulée par atome et par unité de temps sont toutes deux égales à $\sigma \vphi$, où $\sigma >0$ est une constante.
%\begin{center}
%\input{2niveaux.pgf}
%\end{center}

\begin{quest}
\item Définir l'absorption et l'émission stimulée. Justifier que les variations de populations vérifient
\begin{equation}
\left. \frac{\mathrm{d} n_1}{\mathrm{d} t}\right|_{\text{abs}} = -\sigma \vphi n_1
\quad
\text{et}
\quad
\left.\frac{\mathrm{d} n_1}{\mathrm{d} t}\right|_{\text{sti}}  = \sigma \vphi n_2.
\end{equation}
\sol{
Pour ce qui est des unités : $\vphi = \text{(densité de photons)} \times \text{(vitesse d'un photon)}$ = $\text{(nombre de photons)}$/ $\text{(temps)}$ / $\text{(surface)} $ en $\SI{}{m^{-2}.s^{-1}}$. Donc $\sigma$ est une surface (en $\SI{}{m^2}$) : c'est ce qu'on appelle la section efficace de collision.

Tout est dans l'énoncé, il suffit de justifier le signe. Absorption = transition de $1$ vers $2$, donc $n_1$ diminue et $\left. \frac{\mathrm{d} n_1}{\mathrm{d} t}\right|_{\text{abs}}<0$. Émission = transition de $2$ vers $1$, donc $n_1$ augmente et $\left.\frac{\mathrm{d} n_1}{\mathrm{d} t}\right|_{\text{sti}}  >0$.
}
\item L'émission spontanée est négligée. Les photons se déplacent tous à la célérité $c$ dans la même direction. Montrer que la loi d'évolution de la densité de courant de photons est
\begin{equation}
  \label{eq:vphi_gain}
\frac{\mathrm{d} \vphi}{\mathrm{d} t} = \sigma c \Delta n \,\vphi
\end{equation}
avec $\Delta n=n_2-n_1$.
Pour un système à l'équilibre thermodynamique à la température $T$, que vaut le rapport $n_2/n_1$? Que dire alors sur la variation de $\vphi$ au cours du temps ?
\sol{
Soit $n_\gamma$ la densité en photon. On a alors $\vphi = n_\gamma c$, donc
$$
\dd{\vphi}{t} = c \dd{n_\gamma}{t}.
$$
Or chaque transition de $2$ vers $1$ conduisent à la création d'un photon, et donc
$$
\dd{\vphi}{t} = c \dd{n_1}{t} = c \left( \left. \frac{\mathrm{d} n_1}{\mathrm{d} t}\right|_{\text{abs}}  \left.\frac{\mathrm{d} n_1}{\mathrm{d} t}\right|_{\text{sti}}  \right) = \sigma c \vphi (n_2-n_1).
$$

À l'équilibre thermodynamique, le rapport $n_2/n_1$ est donné par le facteur de Boltzmann
$$
\frac{n_2}{n_1} = e^{-(E_2-E_1)/k_B T} < 1.
$$
}

\item Une partie du faisceau traverse les miroirs de sorte que la cavité subit des pertes $\mathrm{d} \vphi|_{\text{cav}} = - K\vphi \mathrm{d}t$. Montrer que l’\Eq{vphi_gain} est modifiée en
\begin{equation}
\frac{\mathrm{d} \vphi}{\mathrm{d} t} = -K \vphi + \sigma c \Delta n \vphi .
\end{equation}
À quelle condition sur $\Delta n$ a-t-on une émission laser?
\sol{
IL suffit d'ajouter le terme $-K\vphi$. Dimension : $[K] = T^{-1}$ est l'inverse du temps caractéristique de transmission de l'intensité lumineuse à travers le miroir de sortie.
On a émission laser si $\dd{\vphi}{t}>0$ soit
$$
-K \vphi + \sigma c \Delta n \vphi > 0 \Rar \Delta n > \frac{K}{\sigma c}.
$$
Cette condition est encore plus contraignante.
}
\end{quest}
Pour qu'un laser puisse fonctionner, il est nécessaire d'inverser les populations atomiques. L'inversion peut être réalisée par pompage optique, qui est modélisé par les termes source $\mathrm{d} n_2|_{\text{pomp}} = \lambda_2 \mathrm{d}t$ et $\mathrm{d} n_1|_{\text{pomp}} = \lambda_1 \mathrm{d}t$. L'émission spontanée vers un état de plus basse énergie est également prise en compte par des termes de pertes $\mathrm{d} n_1|_{\text{spon}} = -\gamma n_1 \mathrm{d}t$ et $\mathrm{d} n_2|_{\text{spon}} = -\gamma n_2 \mathrm{d}t$.
\begin{quest}
  \item Trouver l'équation différentielle vérifiée par $\Delta n$.
% Montrer que $\Delta n$ vérifie l'équation différentielle
% \begin{equation}
% \frac{\mathrm{d} \Delta n}{\mathrm{d} t} = -\gamma(\Delta n-\Delta n_0) - 2\sigma \vphi \Delta n,
% \end{equation}
% où l'on donnera l'expression de la constante $\Delta n_0$. Interpréter.
\sol{
Attention, ici le nombre d'atomes dans $1$ ou $2$, $n_1+n_2$ n'est plus constant en raison du pompage optique ! Ce pompage passe en pratique par la désexcitation d'un niveau d'énergie $E_3>E_2$ vers l'état $E_2$.

On a
$$
\frac{\mathrm{d} n_1}{\mathrm{d} t}  = \left. \frac{\mathrm{d} n_1}{\mathrm{d} t}\right|_{\text{abs}}  + \left. \frac{\mathrm{d} n_1}{\mathrm{d} t}\right|_{\text{sti}}  + \left. \frac{\mathrm{d} n_1}{\mathrm{d} t}\right|_{\text{pomp}}  + \left. \frac{\mathrm{d} n_1}{\mathrm{d} t}\right|_{\text{spon}}  =  -\sigma \vphi n_1 + \sigma \vphi n_2 + \lambda_1 - \gamma n_1 = \sigma \vphi \Delta n + \lambda_1 - \gamma n_1.
$$
De même,
$$
\frac{\mathrm{d} n_2}{\mathrm{d} t}  = \left. \frac{\mathrm{d} n_2}{\mathrm{d} t}\right|_{\text{abs}}  + \left. \frac{\mathrm{d} n_2}{\mathrm{d} t}\right|_{\text{sti}}  + \left. \frac{\mathrm{d} n_1}{\mathrm{d} t}\right|_{\text{pomp}}  + \left. \frac{\mathrm{d} n_2}{\mathrm{d} t}\right|_{\text{spon}}  =  \sigma \vphi n_2 - \sigma \vphi n_1 + \lambda_2 - \gamma n_2 = -\sigma \vphi \Delta n + \lambda_2 - \gamma n_2.
$$
En prenant la différence, on obtient
$$
\dd{\Delta n}{t} = -2\sigma \vphi \Delta n - \gamma \Delta n +  \lambda_2- \lambda_1,
$$
donc $\Delta n_0 = \frac{\lambda_2-\lambda_1}{\gamma}$. En pratique, on cherche à pomper vers l'état $2$ et donc $\lambda_2 \gg \lambda_1$, soit $\Delta n_0 >0$. En régime stationnaire, on doit avoir
$$
\begin{cases}
 0 & = -K \vphi + \sigma c \Delta n \vphi \\
 0 & = -\gamma(\Delta n-\Delta n_0) - 2\sigma \vphi \Delta n \\
\end{cases}
\Rar
\begin{cases}
 \Delta n  & = \frac{K}{\sigma c}\\
  \vphi & = \frac{\gamma}{2 \sigma}(\frac{\Delta n_0}{\Delta n}-1) = \frac{c}{2 K}(\lambda_2-\lambda_1) - \frac{\gamma}{2\sigma}.\\
\end{cases}
$$
Ce flux à un sens (est positif) si $\Delta n_0 > \Delta n$, autrement dit si le pompage $\lambda_2$ est suffisamment grand par rapport aux pertes $K$.
}

\end{quest}
\end{exo}
%%%%%%%%%%%%%%%%%%%%%%%%%%%%%%%%%%%%%%%%%%%%%%%%%%%%%%%%%%%%%%%%%%%%%%%%%%%%%


%%%%%%%%%%%%%%%%%%%%%%%%%%%%%%%%%%%
\begin{exo}[PC][1][colle]{Laser à deux niveaux}
  On considère un milieu amplificateur à deux niveaux $E_1$ (niveau fondamental) et $E_2$ (niveau excité). On appelle $n_1$ et $n_2$ les populations des deux niveaux. On cherche à réaliser une inversion de population en soumettant ce système à un flux de photons accordés sur la transition. 
  % On note $n$ le nombre de ces photons.
% Données : rappel des relations d'Einstein :
% - Probabilité par unité de temps pour l'émission spontanée : $B$
% - Probabilité par unité de temps pour une absorption ou une émission induite (stimulée) : An
\begin{quest}
  \item	
  Établir les équations d'évolution des populations des niveaux 1 et 2.
  \item 
  En déduire la population $n_2$ en régime permanent. 
  % Est-il possible de produire une inversion de population?
\end{quest}
\end{exo}
%%%%%%%%%%%%%%%%%%%%%%%%%%%%%%%%%%%